\usepackage[a4paper,inner=30mm,outer=25mm,top=27mm,bottom=29mm,headsep=8mm]{geometry}
\usepackage{fontspec}
\usepackage{luatexko}
\setmainfont{TeX Gyre Pagella}
\setsansfont{TeX Gyre Heros}
\setmonofont{Menlo}[Scale=MatchLowercase]
\setmainhangulfont{UnBatang.ttf}[
  BoldFont=UnBatangBold.ttf,
  ItalicFont=UnBatang.ttf,
  BoldItalicFont=UnBatangBold.ttf,
  ItalicFeatures={FakeSlant=0.12},
  BoldItalicFeatures={FakeSlant=0.12}
]
\setsanshangulfont{UnDotum.ttf}[BoldFont=UnDotumBold.ttf]
\setmonohangulfont{UnDotum.ttf}[BoldFont=UnDotumBold.ttf]

\usepackage{mathtools,amsthm,bm}
\usepackage{unicode-math}
\setmathfont{TeX Gyre Pagella Math}
\usepackage{microtype}
\usepackage{setspace}
\usepackage{needspace}
\setstretch{1.15}

\usepackage{siunitx}
\usepackage{booktabs,longtable,tabularx,array,multirow}
\usepackage{enumitem}
\usepackage{graphicx}
\usepackage{tikz}
\usetikzlibrary{arrows.meta,positioning,calc,fit,backgrounds,matrix}
\usepackage[most]{tcolorbox}
\usepackage{listings}
\usepackage{xcolor}
\usepackage{xurl}
\usepackage[round,authoryear]{natbib}
\usepackage{imakeidx}
\usepackage{scrlayer-scrpage}
\usepackage{bookmark}
\usepackage{hyperref}
\usepackage[nameinlink,noabbrev]{cleveref}

% luatexko가 문서 시작 시 표준 명칭을 다시 설정하므로 그 뒤에 한국어 명칭을 확정한다.
\AtBeginDocument{%
  \renewcommand{\contentsname}{차례}%
  \renewcommand{\listfigurename}{그림 목차}%
  \renewcommand{\listtablename}{표 목차}%
  \renewcommand{\figurename}{그림}%
  \renewcommand{\tablename}{표}%
  \renewcommand{\bibname}{참고문헌}%
  \renewcommand*{\partformat}{제\ \thepart\ 부}%
  \crefname{equation}{식}{식}%
  \Crefname{equation}{식}{식}%
  \crefname{figure}{그림}{그림}%
  \Crefname{figure}{그림}{그림}%
  \crefname{table}{표}{표}%
  \Crefname{table}{표}{표}%
  \crefname{chapter}{장}{장}%
  \Crefname{chapter}{장}{장}%
  \crefname{section}{절}{절}%
  \Crefname{section}{절}{절}%
}

\definecolor{Ink}{HTML}{1F2933}
\definecolor{Navy}{HTML}{183B56}
\definecolor{Blue}{HTML}{2F6690}
\definecolor{Teal}{HTML}{277D78}
\definecolor{Gold}{HTML}{A66B00}
\definecolor{Brick}{HTML}{A13D3D}
\definecolor{Paper}{HTML}{FAF9F6}
\definecolor{SoftBlue}{HTML}{EEF5FA}
\definecolor{SoftTeal}{HTML}{EEF8F6}
\definecolor{SoftGold}{HTML}{FFF7E7}
\definecolor{SoftRed}{HTML}{FFF1F1}

\hypersetup{
  colorlinks=true,
  linkcolor=Navy,
  citecolor=Teal,
  urlcolor=Blue,
  pdftitle={언어모델의 계산을 해부하다},
  pdfsubject={Decoder-only Transformer와 기계론적 해석가능성},
  pdfauthor={LLM_Study}
}

\setkomafont{disposition}{\sffamily\bfseries}
\setkomafont{chapter}{\sffamily\bfseries\color{Navy}}
\setkomafont{section}{\sffamily\bfseries\color{Navy}}
\setkomafont{subsection}{\sffamily\bfseries\color{Blue}}
\setkomafont{part}{\sffamily\bfseries\color{Navy}}
\setcounter{secnumdepth}{3}
\setcounter{tocdepth}{2}
\KOMAoptions{numbers=noenddot}

\clearpairofpagestyles
\ihead{\headmark}
\ohead{\pagemark}
\automark[section]{chapter}
\KOMAoptions{headsepline=true}

\setlist[itemize]{leftmargin=2em,itemsep=0.25em,topsep=0.4em}
\setlist[enumerate]{leftmargin=2em,itemsep=0.25em,topsep=0.4em}
\renewcommand{\arraystretch}{1.2}

\newcolumntype{Y}{>{\raggedright\arraybackslash}X}
\newcolumntype{R}{>{\raggedleft\arraybackslash}X}

\theoremstyle{definition}
\newtheorem{definition}{정의}[chapter]
\newtheorem{example}[definition]{예제}
\theoremstyle{plain}
\newtheorem{proposition}[definition]{명제}
\newtheorem{theorem}[definition]{정리}
\theoremstyle{remark}
\newtheorem{remark}[definition]{주의}

\newtcolorbox{factbox}{
  breakable, colback=SoftBlue, colframe=Blue, boxrule=0.6pt,
  title={수학적 사실 또는 명시된 구현}, fonttitle=\bfseries\sffamily
}
\newtcolorbox{resultbox}[1]{
  breakable, colback=SoftTeal, colframe=Teal, boxrule=0.6pt,
  title={실험 결과 · #1}, fonttitle=\bfseries\sffamily
}
\newtcolorbox{interpretbox}[1]{
  breakable, colback=SoftGold, colframe=Gold, boxrule=0.6pt,
  title={제한된 해석 · #1}, fonttitle=\bfseries\sffamily
}
\newtcolorbox{limitbox}{
  breakable, colback=SoftRed, colframe=Brick, boxrule=0.6pt,
  title={한계 또는 열린 문제}, fonttitle=\bfseries\sffamily
}
\newtcolorbox{studybox}[1][학습 점검]{
  breakable, colback=Paper, colframe=Ink!55, boxrule=0.5pt,
  title={#1}, fonttitle=\bfseries\sffamily
}

\lstdefinestyle{bookcode}{
  basicstyle=\small\ttfamily,
  keywordstyle=\color{Blue}\bfseries,
  commentstyle=\color{Teal},
  stringstyle=\color{Brick},
  backgroundcolor=\color{Paper},
  frame=single,
  rulecolor=\color{Ink!25},
  breaklines=true,
  showstringspaces=false,
  columns=fullflexible,
  keepspaces=true,
  tabsize=4
}
\lstset{style=bookcode}

\DeclareMathOperator{\softmax}{softmax}
\DeclareMathOperator*{\argmin}{arg\,min}
\DeclareMathOperator{\argmax}{arg\,max}
\DeclareMathOperator{\diag}{diag}
\DeclareMathOperator{\rank}{rank}
\DeclareMathOperator{\tr}{tr}
\DeclareMathOperator{\Var}{Var}
\DeclareMathOperator{\Cov}{Cov}
\DeclareMathOperator{\KL}{KL}
\newcommand{\R}{\mathbb{R}}
\newcommand{\E}{\mathbb{E}}
\newcommand{\given}{\,\middle|\,}
\newcommand{\dd}{\mathop{}\!\mathrm{d}}
\newcommand{\vect}[1]{\bm{#1}}
\newcommand{\mat}[1]{\bm{#1}}
\newcommand{\shape}[1]{\ensuremath{\left[#1\right]}}

\makeindex[intoc,title=찾아보기]
